%!TEX root = ../template.tex
%%%%%%%%%%%%%%%%%%%%%%%%%%%%%%%%%%%%%%%%%%%%%%%%%%%%%%%%%%%%%%%%%%%%
%% abstract-en.tex
%% NOVA thesis document file
%%
%% Abstract in English([^%]*)
%%%%%%%%%%%%%%%%%%%%%%%%%%%%%%%%%%%%%%%%%%%%%%%%%%%%%%%%%%%%%%%%%%%%

\typeout{NT FILE abstract-en.tex}%



Software modeling demands a good understanding of the solution domain, as well as requirements specification. 
The resulting models serve as blueprints for the hardware and software requirements, outlining system architecture, implementation phase, and the overall system structure. 
To automate the creation of such models various methods have been proposed over the years, including approaches based on model-driven development, linguistic rules and patterns, rule-based systems, and statistical or machine learning techniques. 
More recently, Large Language Models (LLMs) have gained prominence. While LLMs often show promising results they commonly lack sufficient contextual awareness, an essential factor in accurately transforming requirement specifications into models such as class diagrams, and output often is inconsistent. 
This dissertation addresses this limitation by introducing an approach that combines prompt engineering techniques, a Retrieval-Augmented Generation system, and a self-refinement mechanism. 
The proposed pipeline provides an automated process for transforming natural language descriptions into UML class diagrams. 
It demonstrates the potential of Generative AI to enhance efficiency and reduce reliance on manual effort in the creation of models.
The approach was evaluated using a combination of qualitative and quantitative methods, including user studies and performance metrics. Results indicate that AI-generated diagrams consistently outperform human-created diagrams in syntactic, semantic, and pragmatic quality, with over 65\% of evaluators unable to distinguish AI outputs from human ones and with over 73\% of evaluators preferring the diagram generated by our system over the human diagram. 
While pragmatic quality remains the most challenging dimension, findings confirm that Generative AI can produce professional-quality diagrams that enhance communication among stakeholders and support early design decisions.


% Palavras-chave do resumo em Inglês
% \begin{keywords}
% Keyword 1, Keyword 2, Keyword 3, Keyword 4, Keyword 5, Keyword 6, Keyword 7, Keyword 8, Keyword 9
% \end{keywords}
\keywords{
  LLM \and
  Requirement modeling \and
  Retrieval-Augmented Generation\and
  UML diagrams\and 
  GAI
}
